% Created 2021-06-07 lun 22:36
% Intended LaTeX compiler: pdflatex
\documentclass[letterpaper, 12pt]{article}
\usepackage[hidelinks]{hyperref}
\usepackage[left=2.5cm, top=2.5cm, right=2.5cm, bottom=2.5cm]{geometry}
\setcounter{secnumdepth}{0}
\usepackage[spanish, es-tabla, es-nodecimaldot]{babel}
\usepackage[utf8]{inputenc}
\usepackage{graphicx}
\usepackage{caption}
\usepackage{subcaption}
\usepackage{physics}
\usepackage{amsmath}
\usepackage{gensymb}
\usepackage{siunitx}
\usepackage{listings}
\usepackage{pdfpages}
\usepackage{cite} % referencias y bibliografía
\usepackage{babelbib}
\usepackage{enumerate}
\addto{\captionsspanish}{\renewcommand*{\listfigurename}{Índice de Figuras}}
\addto{\captionsspanish}{\renewcommand*{\contentsname}{Índice}}
\addto{\captionsspanish}{\renewcommand*{\listtablename}{Índice de Tablas}}
\pagestyle{plain}
\usepackage{eso-pic}
\usepackage[hidelinks]{hyperref}
\providecommand\phantomsection{}
\usepackage{cleveref}
\newcommand{\parg}{\par \hfill \break}
\usepackage{listings}
\usepackage{xcolor}

\definecolor{codegreen}{RGB}{152, 152, 26}
\definecolor{codegray}{RGB}{146, 131, 116}
\definecolor{codepurple}{RGB}{177, 98, 134}
\definecolor{backcolour}{RGB}{251, 241, 199}
\definecolor{codeblue}{RGB}{69, 133, 136}
\lstdefinestyle{mystyle}{
backgroundcolor=\color{backcolour},
commentstyle=\color{codegreen},
keywordstyle=\color{codeblue},
numberstyle=\tiny\color{codegray},
stringstyle=\color{codepurple},
basicstyle=\ttfamily\footnotesize,
breakatwhitespace=false,
breaklines=true,
captionpos=b,
keepspaces=true,
numbers=left,
numbersep=5pt,
showspaces=false,
showstringspaces=false,
showtabs=false,
tabsize=2
}

\lstset{style=mystyle}
\newcommand{\tarea}{Tarea 03: Determinar la matriz inversa de la
transformada de Fourier}
\newcommand{\fecha}{07-junio-2021}
\date{\today}
\title{Tarea 03\\\medskip
\large Diseño y Análisis de algoritmos}
\hypersetup{
 pdfauthor={Marco Antonio Cardoso Moreno},
 pdftitle={Tarea 03},
 pdfkeywords={},
 pdfsubject={},
 pdfcreator={Emacs 27.2 (Org mode 9.5)}, 
 pdflang={Spanish}}
\begin{document}

\begin{titlepage}
\begin{center}
  \AddToShipoutPictureBG*{
    \includegraphics[width=\paperwidth,height=\paperheight]
        {../../portada/portada.pdf}}
        \vspace*{2mm}
        \rule{\textwidth}{1.5pt}\\
        \vspace*{5mm}
        \textbf{\Large INSTITUTO POLITÉCNICO NACIONAL}\\
        \vspace*{8mm}
        {\Large Centro de Investigación en Computación}\\
        \vspace{8mm}
        {\Large Maestría en Ciencias en Ingeniería de Cómputo}\\
        \vspace{19mm}
        {\Large Diseño y Análisis de Algoritmos}\\
        \vspace{49mm}




        \textbf{\large \tarea}\\





        \vspace*{48mm}
        {\large Docente: Dr. Rolando Quintero Téllez}\\
        \vspace*{20mm}
        {\large Alumno: Cardoso Moreno Marco Antonio}\\
        \vspace*{10mm}
        \rule{\textwidth}{1.5pt}\\
        {\large Semestre A21} \hfill {\large Fecha: \fecha}
    \end{center}
\end{titlepage}

\section{Determinar la matriz inversa de la transformada de Fourier}
\label{sec:orge4f464a}
\noindent
La determinación de la matriz inversa de la transformada de Fourier se logra
mediante los algoritmos presentados en el video 5.7 del curso Diseño y Análisis
de Algoritmos.
\parg
Los códigos presentados a continuación fueron desarrollados en lenguaje
\texttt{Python 3.6}, y se compararon los resultados con las funciones desarrolladas
dentro de la librería Numpy.

\subsection{FFT}
\label{sec:org9d3c8b7}
\lstinputlisting[language=python, caption=Función FFT, firstline=12, lastline=27]{./src/dft.py}
\subsection{IFFT}
\label{sec:orga4b9070}
\lstinputlisting[language=python, caption=Función IFFT, firstline=29, lastline=43]{./src/dft.py}
\subsection{Comparación de resultados vs numpy}
\label{sec:orgb62f4ac}
\noindent
Se realizó la comparación de los resultados obtenidos mediante el desarrollo
propio de las funciones \texttt{fft} e \texttt{ifft} contra las funciones \texttt{numpy.fft.fft} y
\texttt{numpy.fft.ifft}, respectivamente.
\parg

La función \texttt{main} donde se ejecuta la comparación de las funciones se muestra a
continuación
\lstinputlisting[language=python, caption=Función Main, firstline=45]{./src/dft.py}

\parg
para la cual los resultados obtenidos fueron:
\begin{center}
\includegraphics[width=.9\linewidth]{./img/resultados_fft.png}
\end{center}
\end{document}
